% GENERATED FILE. Edit canonical lessons, then run npm run generate:books.
% canonical-source-hash: fnv1a64:548c3f08ad0a727a
% canonical-lessons: SA-C147-svami, SA-C147-sevakah, SA-C147-rajni, SA-C147-mantri, SA-C147-senapatih

\chapter{Master, Servant, Queen, Minister, General}
\label{ch:sa-master-servant-queen-minister-general}
\hlchaptermodality{147}

\begin{chapteropening}
\textbf{By the end of this chapter:} \emph{I can say a master, a servant, a queen, a minister and a general.}

Complete the last lesson of chapter 147: I can say a master, a servant, a queen, a minister and a general.
\end{chapteropening}

\section[\texorpdfstring{\sk{स्वामी} (svāmī)}{स्वामी (svāmī)}]{\sk{स्वामी} (svāmī) — a master, an owner}

\label{lesson:SA-C147-svami}



\begin{quote}
Before the new one: say the Sanskrit for a maternal grandfather, then the Sanskrit for a father's brother.
\end{quote}

\subsection*{You'll want to know: \sk{स्वामी}}
\textbf{\sk{स्वामी}} — \emph{svāmī} — \textquotedblleft{}a master, an owner\textquotedblright{}.

\begin{grammarlens}[title={What you've built}]
One more word for this chapter.
\end{grammarlens}

\subsection*{Your turn}
\begin{itemize}
\raggedright
  \item \emph{Say it:} \emph{svāmī}
  \item \emph{Say it:} \emph{svāmī}, once more
  \item \emph{Say it:} say \emph{svāmī}
  \item \emph{Recall:} say the Sanskrit for a maternal grandfather, then the Sanskrit for a father's brother, then say \emph{svāmī} again
\end{itemize}

\begin{tcolorbox}[breakable,colback=teal!4,colframe=teal!35!black,title={Before you move on}]
What does \sk{स्वामी} mean? (\textquotedblleft{}A master, an owner\textquotedblright{}.) Say it once more.
\end{tcolorbox}
\section[\texorpdfstring{\sk{सेवकः} (sevakaḥ)}{सेवकः (sevakaḥ)}]{\sk{सेवकः} (sevakaḥ) — a servant, an attendant}

\label{lesson:SA-C147-sevakah}



\begin{quote}
Before the new one: say the Sanskrit for a father's brother, then the Sanskrit for a master.
\end{quote}

\subsection*{You'll want to know: \sk{सेवकः}}
\textbf{\sk{सेवकः}} — \emph{sevakaḥ} — \textquotedblleft{}a servant, an attendant\textquotedblright{}.

\begin{grammarlens}[title={What you've built}]
One more word for this chapter.
\end{grammarlens}

\subsection*{Your turn}
\begin{itemize}
\raggedright
  \item \emph{Say it:} \emph{sevakaḥ}
  \item \emph{Say it:} \emph{sevakaḥ}, once more
  \item \emph{Say it:} say \emph{sevakaḥ}
  \item \emph{Recall:} say the Sanskrit for a father's brother, then the Sanskrit for a master, then say \emph{sevakaḥ} again
\end{itemize}

\begin{tcolorbox}[breakable,colback=teal!4,colframe=teal!35!black,title={Before you move on}]
What does \sk{सेवकः} mean? (\textquotedblleft{}A servant, an attendant\textquotedblright{}.) Say it once more.
\end{tcolorbox}
\section[\texorpdfstring{\sk{राज्ञी} (rājñī)}{राज्ञी (rājñī)}]{\sk{राज्ञी} (rājñī) — a queen}

\label{lesson:SA-C147-rajni}



\begin{quote}
Before the new one: say the Sanskrit for a master, then the Sanskrit for a servant.
\end{quote}

\subsection*{You'll want to know: \sk{राज्ञी}}
\textbf{\sk{राज्ञी}} — \emph{rājñī} — \textquotedblleft{}a queen\textquotedblright{}.

\begin{grammarlens}[title={What you've built}]
One more word for this chapter.
\end{grammarlens}

\subsection*{Your turn}
\begin{itemize}
\raggedright
  \item \emph{Say it:} \emph{rājñī}
  \item \emph{Say it:} \emph{rājñī}, once more
  \item \emph{Say it:} say \emph{rājñī}
  \item \emph{Recall:} say the Sanskrit for a master, then the Sanskrit for a servant, then say \emph{rājñī} again
\end{itemize}

\begin{tcolorbox}[breakable,colback=teal!4,colframe=teal!35!black,title={Before you move on}]
What does \sk{राज्ञी} mean? (\textquotedblleft{}A queen\textquotedblright{}.) Say it once more.
\end{tcolorbox}
\section[\texorpdfstring{\sk{मन्त्री} (mantrī)}{मन्त्री (mantrī)}]{\sk{मन्त्री} (mantrī) — a minister}

\label{lesson:SA-C147-mantri}



\begin{quote}
Before the new one: say the Sanskrit for a servant, then the Sanskrit for a queen.
\end{quote}

\subsection*{You'll want to know: \sk{मन्त्री}}
\textbf{\sk{मन्त्री}} — \emph{mantrī} — \textquotedblleft{}a minister\textquotedblright{}.

\begin{grammarlens}[title={What you've built}]
One more word for this chapter.
\end{grammarlens}

\subsection*{Your turn}
\begin{itemize}
\raggedright
  \item \emph{Say it:} \emph{mantrī}
  \item \emph{Say it:} \emph{mantrī}, once more
  \item \emph{Say it:} say \emph{mantrī}
  \item \emph{Recall:} say the Sanskrit for a servant, then the Sanskrit for a queen, then say \emph{mantrī} again
\end{itemize}

\begin{tcolorbox}[breakable,colback=teal!4,colframe=teal!35!black,title={Before you move on}]
What does \sk{मन्त्री} mean? (\textquotedblleft{}A minister\textquotedblright{}.) Say it once more.
\end{tcolorbox}
\section[\texorpdfstring{\sk{सेनापतिः} (senāpatiḥ)}{सेनापतिः (senāpatiḥ)}]{\sk{सेनापतिः} (senāpatiḥ) — a general}

\label{lesson:SA-C147-senapatih}



\begin{quote}
Before the new one: say the Sanskrit for a queen, then the Sanskrit for a minister.
\end{quote}

\subsection*{You'll want to know: \sk{सेनापतिः}}
\textbf{\sk{सेनापतिः}} — \emph{senāpatiḥ} — \textquotedblleft{}a general\textquotedblright{}.

\begin{grammarlens}[title={What you've built}]
One more word for this chapter.
\end{grammarlens}

\subsection*{Your turn}
\begin{itemize}
\raggedright
  \item \emph{Say it:} \emph{senāpatiḥ}
  \item \emph{Say it:} \emph{senāpatiḥ}, once more
  \item \emph{Say it:} say \emph{senāpatiḥ}
  \item \emph{Recall:} say the Sanskrit for a queen, then the Sanskrit for a minister, then say \emph{senāpatiḥ} again
\end{itemize}

\begin{tcolorbox}[breakable,colback=teal!4,colframe=teal!35!black,title={Before you move on}]
What does \sk{सेनापतिः} mean? (\textquotedblleft{}A general\textquotedblright{}.) Say it once more.
\end{tcolorbox}
